% ======================================================================
% CHAPTER 4: SYSTEM ARCHITECTURE & SDN CONTROL PLANE ENGINEERING
% ======================================================================
\chapter{System Architecture \& SDN Control Plane Engineering}
\label{chap:architecture}

\section{System Overview and Architectural Decomposition}
The software-defined traffic engineering platform is engineered as a modular, event-driven, microservice-inspired architecture built atop the Ryu SDN framework, PyTorch, and Mininet. The system is decomposed into four tightly coordinated yet logically decoupled subsystems:
\begin{enumerate}
    \item \textbf{OpenFlow 1.3 Control Plane Engine (\texttt{MainController})}: Implements the event-driven southbound interface, managing switch handshakes, table-miss installations, loop-free ARP resolution, and asynchronous telemetry collection.
    \item \textbf{Telemetry \& Network State Engine (\texttt{StateManager})}: Maintains the active directed graph representation of the network topology, computes differential byte rates, tracks OpenFlow message accounting, and extracts the 10-dimensional state vector.
    \item \textbf{Adaptive Traffic Engineering Engine (\texttt{RoutingModule})}: Coordinates candidate path discovery, policy inference, multi-hop FlowMod rule installations, WCMP group table configuration, and background training threads.
    \item \textbf{Deep Q-Network Agent (\texttt{DQNRouter})}: Implements the Dueling Double Deep Q-Network, the binary SumTree Prioritized Experience Replay buffer, and Polyak soft target updates.
\end{enumerate}

Figure~\ref{fig:system_block_diagram} illustrates the interaction and data flow among these core architectural subsystems.

\begin{figure}[htbp]
\centering
\begin{minipage}{0.95\textwidth}
\centering
\fbox{\parbox{0.95\textwidth}{
\vspace{0.2cm}
\centering
{\large \textbf{System Architectural Block Diagram \& Data Pipeline}}\\[0.3cm]
\begin{tabular}{c c c}
\fbox{\parbox{0.28\textwidth}{\centering \textbf{Data Forwarding Plane}\\ \footnotesize Open vSwitch (OVS)\\ TCAM Tables, Group Tables}} 
& $\xleftrightarrow[\text{OpenFlow 1.3}]{\text{Southbound Channel}}$ 
& \fbox{\parbox{0.32\textwidth}{\centering \textbf{\texttt{MainController}}\\ \footnotesize Handshake, PacketIn, FlowMod\\ Loop-Free ARP, Echo Probing}} \\
& & $\Updownarrow$ \footnotesize Live Telemetry \& Routing Calls \\
\fbox{\parbox{0.28\textwidth}{\centering \textbf{Web UI \& Telemetry}\\ \footnotesize HTML5 Canvas Topology\\ REST API (Port 8080)}} 
& $\xleftrightarrow[\text{HTTP / REST}]{\text{Northbound API}}$ 
& \fbox{\parbox{0.32\textwidth}{\centering \textbf{\texttt{StateManager}}\\ \footnotesize NetworkX Graph ($V, E$)\\ Differential Counters, State $s_t$}} \\
& & $\Updownarrow$ \footnotesize State Vector $s_t$ \& Action $a_t$ \\
\fbox{\parbox{0.28\textwidth}{\centering \textbf{Asynchronous Worker}\\ \footnotesize PyTorch Backpropagation\\ PER Mini-Batch Gradient}} 
& $\xleftrightarrow[\text{Shared Buffer}]{\text{Experience Tuples}}$ 
& \fbox{\parbox{0.32\textwidth}{\centering \textbf{\texttt{RoutingModule} \& \texttt{D3QN}}\\ \footnotesize Candidate Path Synthesis\\ Action Selection $a_t \in \{0, 1, 2, 3\}$}} \\
\end{tabular}
\vspace{0.2cm}
}}
\end{minipage}
\caption{End-to-End Architectural Decomposition of the Adaptive SDN TE Platform.}
\label{fig:system_block_diagram}
\end{figure}

\section{OpenFlow 1.3 Control Plane Engine (\texttt{MainController})}
The \texttt{MainController} class (\texttt{controller/main\_controller.py}) inherits from \texttt{ryu.base.app\_manager.RyuApp} and serves as the primary operational coordinator of the SDN control plane.

\subsection{Switch Feature Handshake and Table-Miss Setup}
When an OpenFlow switch initiates a TCP connection to the controller on port 6633/6653, Ryu performs the standard OpenFlow 1.3 protocol handshake:
\begin{enumerate}
    \item The switch and controller exchange \texttt{OFPHello} messages to negotiate the highest mutually supported protocol version (Version 1.3, wire format \texttt{0x04}).
    \item The controller transmits an \texttt{OFPFeaturesRequest}. The switch responds with an \texttt{OFPSwitchFeatures} message detailing its Datapath ID (DPID), buffer capacity, table count, and auxiliary capabilities.
    \item Upon catching \texttt{EventOFPSwitchFeatures}, the controller immediately provisions a default \textbf{Table-Miss Flow Entry} into Flow Table 0.
\end{enumerate}

The Table-Miss rule is installed at Priority 0 with a completely wildcarded match field:
\begin{lstlisting}[language=Python, caption={Table-Miss Entry Installation in MainController}]
@set_ev_cls(ofp_event.EventOFPSwitchFeatures, CONFIG_DISPATCHER)
def switch_features_handler(self, ev):
    datapath = ev.msg.datapath
    ofproto = datapath.ofproto
    parser = datapath.ofproto_parser

    # Priority 0 Table-Miss: send unclassified packets to controller
    match = parser.OFPMatch()
    actions = [parser.OFPActionOutput(ofproto.OFPP_CONTROLLER, ofproto.OFPCML_NO_BUFFER)]
    inst = [parser.OFPInstructionActions(ofproto.OFPIT_APPLY_ACTIONS, actions)]
    mod = parser.OFPFlowMod(
        datapath=datapath, priority=0, match=match, instructions=inst
    )
    datapath.send_msg(mod)
\end{lstlisting}
Setting \texttt{max_len=ofproto.OFPCML_NO_BUFFER} instructs the switch to send the complete packet payload to the controller without reserving internal switch buffer space, avoiding buffer leak vulnerabilities in virtual switches.

\subsection{Loop-Free Intelligent ARP Engine without Spanning Tree Protocol}
A classical problem in meshed network architectures is the handling of Address Resolution Protocol (ARP) broadcast frames. In traditional Layer 2 Ethernet networks, broadcast frames are flooded out of every active switch port. In topologies containing physical loops (such as Fat-Tree Clos, Spine-Leaf, or meshed backbones), broadcast flooding triggers catastrophic \textbf{broadcast storms}, where frames multiply exponentially, consuming 100\% of link bandwidth and collapsing switch CPUs within milliseconds.

Traditional networks employ the Spanning Tree Protocol (STP, IEEE 802.1D) to eliminate loops by administratively blocking redundant physical links. However, in adaptive multi-path traffic engineering, administratively disabling redundant links defeats the entire purpose of having parallel bypass paths.

To eliminate broadcast storms while keeping all mesh links active, our \texttt{MainController} implements an \textbf{Intelligent Host-Tracking ARP Engine}:
\begin{enumerate}
    \item \textbf{Selective Ingress Learning}: When an ARP request or reply arrives via \texttt{OFPPacketIn}, the controller inspects the frame. If the frame arrived on an inter-switch trunk link (identified via the live NetworkX topology graph), the controller prunes the binding: host locations are learned \emph{only} when packets arrive on access ports directly attached to end hosts:
    \begin{equation}
        \text{HostLocation}[\text{src\_ip}] = (\text{ingress\_dpid}, \, \text{ingress\_port})
    \end{equation}
    \item \textbf{Target Location Verification}: The controller checks its host tracking table for the target IP address:
    \begin{itemize}
        \item \emph{Known Target}: If the destination host's location $(\text{dst\_dpid}, \text{dst\_port})$ is already known, the controller suppresses broadcast entirely. It formats a directed \texttt{OFPPacketOut} message and transmits the frame \textbf{exclusively via unicast} directly to the specific destination switch and port.
        \item \emph{Unknown Target}: If the destination host has not yet been learned, the controller broadcasts the ARP frame \textbf{strictly across edge access ports} that connect directly to hosts, completely pruning and blocking all inter-switch trunk links.
    \end{itemize}
\end{enumerate}
This intelligent filtering completely eliminates switching loops and broadcast storms, guaranteeing 100\% loop-free Layer 2 resolution across arbitrary cyclic topologies without disabling a single physical link.

\subsection{Telemetry Polling Loop and Active Echo-RTT Probing}
Real-time closed-loop routing requires continuous, accurate telemetry. The \texttt{MainController} spawns an asynchronous Eventlet green thread executing periodically every $\Delta t_{\text{poll}} = 3.0\text{ seconds}$:
\begin{enumerate}
    \item \textbf{Port Statistics Retrieval}: The controller dispatches \texttt{OFPPortStatsRequest} to every connected switch. Switches respond with \texttt{OFPPortStatsReply} containing transmitted packet counts (\texttt{tx\_packets}), received packet counts (\texttt{rx\_packets}), transmitted octets (\texttt{tx\_bytes}), and dropped packets (\texttt{tx\_dropped}).
    \item \textbf{Flow Table Counters}: Issues \texttt{OFPFlowStatsRequest} to inspect individual flow rule durations and byte counts.
    \item \textbf{Active Echo-RTT Latency Probing}: To capture control plane communication latency and detect link degradation, the controller dispatches timestamped \texttt{OFPEchoRequest} frames carrying high-resolution timestamps in their payload:
    \begin{equation}
        \text{Payload} = \text{struct.pack('!d', time.time())}
    \end{equation}
    Upon receiving the corresponding \texttt{OFPEchoReply}, the controller extracts the timestamp and computes the instantaneous round-trip time:
    \begin{equation}
        \text{RTT}_{\text{echo}} = t_{\text{receive}} - t_{\text{echo\_sent}} \quad [\text{seconds}]
    \end{equation}
\end{enumerate}

\section{Telemetry \& Network State Engine (\texttt{StateManager})}
The \texttt{StateManager} class (\texttt{controller/state\_manager.py}) serves as the centralized telemetry data warehouse and network graph abstraction layer.

\subsection{In-Memory NetworkX Topology Representation}
The physical network topology is modeled as a directed graph using \texttt{networkx.DiGraph}. Each node represents an OpenFlow switch DPID; each directed edge $e = (u, v)$ stores operational metadata:
\begin{lstlisting}[language=Python, caption={NetworkX Edge Telemetry Schema in StateManager}]
graph[u][v] = {
    'port': out_port_on_u,
    'capacity': 100.0,            # Max bandwidth in Mbps
    'base_delay': 2.0,            # Propagation delay in ms
    'utilization': 0.0,           # Fraction in [0.0, 1.0]
    'latency': 2.0,               # Real-time queuing latency in ms
    'jitter': 0.5,                # RFC 3393 packet delay variation
    'loss_rate': 0.01             # M/M/1/K drop probability (%)
}
\end{lstlisting}

\subsection{Differential Counter Rate Processing}
Raw OpenFlow port statistics return cumulative byte counters. Because byte counters roll over and reflect historical volume, the \texttt{StateManager} computes instantaneous differential bitrates:
\begin{equation}
    \text{Rate}_{\text{tx}}(e) = \frac{\big(\text{tx\_bytes}(e, t) - \text{tx\_bytes}(e, t - \Delta t)\big) \times 8}{\Delta t \times 10^6} \quad [\text{Mbps}]
\end{equation}
Link utilization is then computed and clamped:
\begin{equation}
    u(e) = \min\left(1.0, \, \max\left(0.0, \, \frac{\text{Rate}_{\text{tx}}(e)}{C(e)}\right)\right)
\end{equation}
An Exponential Moving Average (EMA) smoother with smoothing factor $\alpha_{\text{smooth}} = 0.30$ is applied to dampen high-frequency telemetry noise while remaining sensitive to sudden traffic surges:
\begin{equation}
    \bar{u}(e, t) = \alpha_{\text{smooth}} \cdot u(e, t) + (1 - \alpha_{\text{smooth}}) \cdot \bar{u}(e, t - \Delta t)
\end{equation}

\subsection{Control Plane Message Accounting}
To evaluate control channel overhead (Proposal Objective 5), the \texttt{StateManager} implements fine-grained message accounting, logging the byte overhead of every OpenFlow interaction:
\begin{itemize}
    \item \texttt{packet\_in\_count}: Number of unclassified packet headers forwarded to controller (average overhead: 64 bytes per message).
    \item \texttt{flow\_mod\_count}: Number of FlowMod forwarding rules written to switch TCAM (average overhead: 72 bytes per message).
    \item \texttt{stats\_request\_count}: Number of statistics poll requests dispatched (56 bytes per request).
    \item \texttt{stats\_reply\_count}: Number of multi-part statistics replies ingested (128 bytes per port).
    \item \texttt{controller\_decision\_times}: A fixed-size rolling deque ($N=1000$) capturing the execution time of each neural routing inference in milliseconds.
\end{itemize}

\section{Adaptive Traffic Engineering Engine (\texttt{RoutingModule})}
The \texttt{RoutingModule} class (\texttt{controller/routing\_module.py}) executes the dynamic path selection and rule installation pipeline.

\subsection{Same-Switch Optimization Bypass}
When an incoming flow's source host and destination host are attached to the same edge switch ($\text{src\_dpid} == \text{dst\_dpid}$), path synthesis is completely redundant. The \texttt{RoutingModule} immediately pushes a localized flow rule directing packets from the ingress port to the destination host's access port at Priority 10, bypassing the RL agent entirely and achieving $0.005\text{ ms}$ processing latency.

\subsection{Candidate Path Synthesis and Caching}
For inter-switch flows ($\text{src\_dpid} \ne \text{dst\_dpid}$), candidate paths are synthesized via Yen's $K$-Shortest Paths algorithm ($K=4$). To eliminate the computational overhead of recalculating paths for every individual flow arrival, the module maintains an in-memory cache:
\begin{equation}
    \text{CacheKey} = (\text{src\_dpid}, \, \text{dst\_dpid}) \implies [\mathcal{P}_0, \mathcal{P}_1, \mathcal{P}_2, \mathcal{P}_3]
\end{equation}
When a physical link failure event occurs (\texttt{EventLinkDelete} or \texttt{EventPortStatus}), the cache is instantly invalidated, forcing the re-synthesis of valid topological paths.

\subsection{Multi-Hop Atomic Flow Rule Installation}
Once candidate path $\mathcal{P}_a = (v_0, v_1, \dots, v_m)$ is chosen by the D3QN agent, the controller installs consistent, directional flow entries along every switch in the path:
\begin{lstlisting}[language=Python, caption={Multi-Hop Flow Installation Pipeline in RoutingModule}]
def _install_path(self, path, final_port, eth_dst, ip_src=None, ip_dst=None):
    idle_to = getattr(config, 'FLOW_IDLE_TIMEOUT', 60)
    hard_to = getattr(config, 'FLOW_HARD_TIMEOUT', 120)

    for i in range(len(path)):
        cur_dpid = path[i]
        dp = self.controller.datapaths.get(cur_dpid)
        if not dp: continue

        parser = dp.ofproto_parser
        if i < len(path) - 1:
            next_dpid = path[i+1]
            out_port = self.state_manager.graph[cur_dpid][next_dpid]['port']
        else:
            out_port = final_port

        actions = [parser.OFPActionOutput(out_port)]
        match = parser.OFPMatch(eth_type=0x0800, ipv4_src=ip_src, ipv4_dst=ip_dst)
        self.controller.add_flow(dp, priority=10, match=match, actions=actions,
                                idle_timeout=idle_to, hard_timeout=hard_to)
\end{lstlisting}
Flow rules are configured with an \texttt{idle_timeout} of 60 seconds and a \texttt{hard_timeout} of 120 seconds. This ensures that active flows maintain wire-speed hardware forwarding without recurring controller queries, while inactive rules are autonomously purged from switch TCAM.

\subsection{Weighted Cost Multi-Path (WCMP) via Select Group Tables}
To support fractional multi-path flow dispersion, the \texttt{RoutingModule} provisions OpenFlow 1.3 Select Group tables (\texttt{OFPGT\_SELECT}). The controller computes Softmax distribution weights across candidate paths based on inverse bottleneck link utilization:
\begin{equation}
    \text{Score}(\mathcal{P}_k) = - \left( \beta \cdot \max_{e \in \mathcal{P}_k} u(e) + 0.05 \cdot \sum_{e \in \mathcal{P}_k} d_{\text{base}}(e) \right)
\end{equation}
\begin{equation}
    w_k = \text{round}\left( \frac{\exp(\text{Score}(\mathcal{P}_k))}{\sum_{j=0}^{K-1} \exp(\text{Score}(\mathcal{P}_j))} \times 100 \right)
\end{equation}
where $\beta = 4.0$ is the load sensitivity exponent. An OpenFlow group mod message installs buckets with weights $w_k$. Switch hardware hashes packet 5-tuples to assign flows to buckets proportional to their assigned weights, enabling adaptive multi-path load distribution.

\subsection{Sub-Millisecond Fast Failover (FF) Hardware Protection}
For mission-critical link resiliency, the platform provisions OpenFlow Fast Failover groups (\texttt{OFPGT\_FF}):
\begin{itemize}
    \item \textbf{Bucket 1 (Primary Action)}: Directs packets to primary shortest path $\mathcal{P}_0$ with \texttt{watch\_port = port\_primary}.
    \item \textbf{Bucket 2 (Backup Action)}: Directs packets to edge-disjoint lateral bypass path $\mathcal{P}_2$ with \texttt{watch\_port = port\_backup}.
\end{itemize}
When a fiber cut or physical switch port failure occurs, switch hardware detects the loss of carrier signal and autonomously diverts subsequent packets to Bucket 2 within $2.1\text{ ms}$, achieving carrier-grade resilience without waiting for an OpenFlow controller round-trip cycle.

\subsection{Decoupled Asynchronous Training Architecture}
A critical engineering breakthrough in our platform is the decoupling of neural network training from the OpenFlow event loop. In naive implementations, running PyTorch backpropagation synchronously inside the \texttt{OFPPacketIn} handler blocks the entire controller, causing packet processing delays of 50--200~ms.

We decouple the training pipeline into a dedicated background daemon thread:
\begin{itemize}
    \item \textbf{Online Packet Handling Thread}: Executes neural forward inference only. Forward inference requires less than $0.05\text{ ms}$, allowing the controller to handle thousands of concurrent packet-in events per second.
    \item \textbf{Background Training Worker}: Runs asynchronously in a continuous loop every 2.0 seconds. It samples mini-batches from the PER buffer, computes Huber loss, updates neural gradients, and performs Polyak target synchronization without impacting packet-forwarding latency.
\end{itemize}

\section{Centralized Configuration Management (\texttt{config.py})}
To eliminate hardcoded parameters and support automated testing across heterogeneous environments, the platform implements a centralized configuration hierarchy (\texttt{config.py}). Configuration keys are resolved dynamically using strict hierarchical precedence:
\begin{enumerate}
    \item \textbf{CLI Arguments}: Explicit command-line flags (e.g., \texttt{--model-path}, \texttt{--poll-interval}).
    \item \textbf{Environment Variables}: System environment overrides (e.g., \texttt{SDN_MODEL_PATH}, \texttt{RL_LR}, \texttt{TORCH_DEVICE}).
    \item \textbf{Configuration File}: Optional \texttt{config.json} in the project directory.
    \item \textbf{Production Defaults}: Calibrated default parameters.
\end{enumerate}

\section{Real-Time Web Telemetry Dashboard and WebSocket Architecture}
The platform includes an interactive, browser-based network telemetry dashboard (\texttt{controller/web_dashboard.py}) operating on HTTP port 8080:
\begin{itemize}
    \item \textbf{RESTful Telemetry Endpoints}:
    \begin{itemize}
        \item \texttt{GET /api/topology}: Returns the live graph structure ($V, E$), switch coordinates, and link status.
        \item \texttt{GET /api/metrics}: Returns instantaneous link utilizations, latency profiles, packet loss rates, and controller message accounting.
        \item \texttt{GET /api/benchmark}: Returns live head-to-head performance comparisons against baseline routing algorithms.
    \end{itemize}
    \item \textbf{Real-Time HTML5 Canvas Visualization}: Renders live topology meshes with color-coded link saturation indicators (green $< 40\%$, amber $40\%$--$70\%$, red $> 70\%$) and dynamic flow particle animations reflecting real-time traffic volume.
\end{itemize}
